% Folder 70, batch 5 (archivists' pages 81-100).
% Pass: Opus 5.5 (claude-opus-5-5), 2026-09-29 — first pass, unchecked against the pages by a human.
\documentclass[11pt,a4paper]{article}
\input{../preamble/grothendieck}

\folder{70}
\batch{5}
\pages{81}{100}
\foldertitle{Réalisations géométriques de structures combinatoires (n-polyèdres, n-hyperpolyèdres [2-polyèdres réguliers]…) (Vieilles rédactions) : notes manuscrites (s.d.).}
\dating{[à partir de 1976]}
\watermark{Édition de démonstration}

\begin{document}

\page{82}
\note{La page s'ouvre sans titre sur un drapeau et trois générateurs ; l'argument vient vraisemblablement de pages antérieures à ce lot.}

\drawing{en haut à gauche, un polyèdre (icosaèdre ?) dont une face triangulaire porte les sommets $s_0, s_1, s_2$ ; au-dessous, un triangle $s_0 s_1 s_2$ avec sa hauteur issue de $s_2$ et les médianes, les régions marquées $\sigma_0, \sigma_1, \sigma_2$ ; en tirets, les sommets voisins $s_3$ (sous $s_0 s_1$), $s'_3$ (à gauche) et $s''_3$ (à droite)}

\[ R = (f_0, f_1, f_2) \qquad
\begin{cases}
\sigma_0 f_0 \neq f_0, & \sigma_0 f_1 = f_1, \quad \sigma_0 f_2 = f_2 \\
\sigma_1 f_0 = f_0, & \sigma_1 f_1 \neq f_1, \quad \sigma_1 f_2 = f_2 \\
\sigma_2 f_0 = f_0, & \sigma_2 f_1 = f_1
\end{cases} \]

$s_0$, $s_1 = \sigma_0 s_0$, $s_2 = \sigma_1 s_1$, $s_3 = \sigma_2 s_2$ ; \quad $f_1 = \mathrm{Dr}(s_0, s_1)$, \quad $f_2 = \mathrm{Pl}(s_0, s_1, s_2)$.

\[ \sigma_0 : \begin{cases} s_0 \mapsto s_1 \\ s_1 \mapsto s_0 \\ s_2 \mapsto s_2 \\ s_3 \mapsto s_3 \end{cases}
\qquad
\sigma_1 : \begin{cases} s_0 \mapsto s_0 \\ s_1 \mapsto s_2 \\ s_2 \mapsto s_1 \\ s_3 \mapsto s'_3 \overset{?}{=} \struck{s''_3} \end{cases}
\qquad
\sigma_2 : \begin{cases} s_0 \mapsto s_0 \\ s_1 \mapsto s_1 \\ s_2 \mapsto s_3 \\ s_3 \mapsto s_2 \end{cases} \]

\[ \begin{cases} e_1 = s_1 - s_0 \\ e_2 = s_2 - s_0 \\ e_3 = s_3 - s_0 \end{cases} \]

\[ s'_3 = s_0 + \lambda e_1 + \mu e_2 + \nu e_3 \quad \bigl(= s_0 + \lambda e_1 - \nu\lambda e_2 + \nu e_3\bigr) \]
\note{la parenthèse est lue avec doute : \uncertain{$-\nu\lambda e_2$}.}

\[ \sigma_0 : \begin{cases} e_1 \mapsto -e_1 \\ e_2 \mapsto -e_1 + e_2 \\ e_3 \mapsto -e_1 + e_3 \end{cases}
\qquad \struck{\uncertain{\text{matrice}}} \quad
\sigma_0 = \begin{pmatrix} -1 & -1 & -1 & \ill{} \\ 0 & 1 & 0 & 0 \\ 0 & 0 & 1 & 0 \\ 0 & 0 & 0 & 1 \end{pmatrix} \]
\note{le quatrième coefficient de la première ligne est couvert d'une tache d'encre ; la dernière ligne est écrite en surcharge.}

\[ \sigma_1 : \begin{cases} e_1 \mapsto \struck{e_0}\ e_2 \\ e_2 \mapsto e_1 \\ e_3 \mapsto \lambda e_1 + \mu e_2 + \nu e_3 \quad (= \lambda e_1 - \lambda e_2 + e_3) \end{cases}
\qquad \struck{\uncertain{\text{matrice}}} \quad
\sigma_1 = \begin{pmatrix} 0 & 1 & \lambda & 0 \\ 1 & 0 & \mu & 0 \\ 0 & 0 & \nu & 0 \\ 0 & 0 & 0 & 1 \end{pmatrix} \]

\[ \sigma_2 : \begin{cases} e_1 \mapsto e_1 \\ e_2 \mapsto e_3 \\ e_3 \mapsto e_2 \end{cases}
\qquad
\sigma_2 = \begin{pmatrix} 1 & 0 & 0 & 0 \\ 0 & 0 & 1 & 0 \\ 0 & 1 & 0 & 0 \\ 0 & 0 & 0 & 1 \end{pmatrix} \]

\page{83}
Écrivons les relations
\[ \sigma_1^2 = \mathrm{id} \]
(les relations $\sigma_0^2 = \sigma_2^2 = \mathrm{id}$ et $\sigma_0\sigma_2 = \sigma_2\sigma_0$ étant évidentes). Il suffit d'écrire
\[ \sigma_1^2(s_3) = s_3 \quad \text{i.e.} \quad \sigma_1(s'_3) = s_3 \quad \text{or} \]
\note{l'exposant de $s_3$ dans le premier membre de droite ressemble à un 4 : \uncertain{$s_3$}.}
\begin{align*}
\sigma_1(s'_3) &= \sigma_1 s_0 + \struck{\sigma_1 e_3 = s_0 +} \lambda\sigma_1 e_1 + \mu\sigma_1 e_2 + \nu\sigma_1 e_3 \\
&= s_0 + \lambda e_2 + \mu e_1 + \nu(\lambda e_1 + \mu e_2 + \nu e_3) \\
&= s_0 + (\mu + \nu\lambda) e_1 + (\lambda + \nu\mu) e_2 + \nu^2 e_3 .
\end{align*}

Les conditions sont
\[ \begin{cases} \nu^2 = 1 \quad \struck{\mu + \nu\lambda} & (1^\circ) \\ \mu = -\nu\lambda & (2^\circ) \\ \lambda = -\nu\mu \end{cases} \]
ou encore $\lambda = \nu^2\lambda$ (\uncertain{conséquence de $2^\circ$}), qui résulte de $1^\circ$.

Sur un corps, ces conditions se décomposent en deux possibilités
\[ \begin{cases} \nu = 1, & \mu = -\lambda \\ \nu = -1, & \mu = \lambda \end{cases} \]
Notons que
\[ \det \sigma_1 = \struck{\det} -\nu \]
donc si on veut $\det \sigma_1 = -1$ (i.e.\ si on veut $\sigma_1$ une réflexion) on doit avoir $\nu = 1$, $\mu = -\lambda$.

\marginal{écrit en oblique dans la marge gauche, à hauteur des deux possibilités : « On \ill{} \ill{} \ill{} \ill{} $\nu = 1$, \ill{} $\sigma_0(s'_3) = \ill{}$, \ill{} $s''_3 = \ill{}$ \ldots\ i.e.\ $\sigma_1\sigma_0\sigma_1\sigma_0(s_3)$ \ldots »}
\note{la note marginale oblique n'est lue qu'en fragments ; elle semble exprimer $s''_3$ au moyen de $\sigma_0$ et $\sigma_1$.}

On aura donc
\[ s'_3 - s_3 = \struck{\lambda e_1 + \ldots} (s'_3 - s_0) - (s_3 - s_0) = \lambda(e_2 - e_1) = -\lambda(s_2 - s_1) \]
\note{sous $(s'_3 - s_0)$ il écrit $u(e_3)$, sous $(s_3 - s_0)$ il écrit $e_3$. Le signe de $\lambda(e_2 - e_1)$ est lu tel quel : d'après la page 82 on attendrait $\lambda(e_1 - e_2)$, ce qui donne bien $-\lambda(s_2 - s_1)$.}

On posera
\[ (1) \qquad -\lambda = (1 + \alpha), \qquad \lambda = -(1+\alpha) = \alpha' \]
de sorte qu'on aura
\[ (2) \qquad s'_3 - s_3 = (1 + \alpha)(s_2 - s_1) \]

\page{84}
et la matrice de $\sigma_1$ sera
\[ (3) \qquad \sigma_1 = \begin{pmatrix} 0 & 1 & -(1+\alpha) & 0 \\ 1 & 0 & 1+\alpha & 0 \\ 0 & 0 & 1 & 0 \\ 0 & 0 & 0 & 1 \end{pmatrix}
= \begin{pmatrix} 0 & 1 & \uncertain{-\alpha'} & 0 \\ 1 & 0 & -\alpha' & 0 \\ 0 & 0 & 1 & 0 \\ 0 & 0 & 0 & 1 \end{pmatrix} \]
\note{avec $\alpha' = -(1+\alpha)$ (page 83), la seconde matrice devrait porter $\alpha'$ puis $-\alpha'$ dans la troisième colonne ; le premier signe est lu avec doute.}

Posons
\[ (4) \qquad \begin{cases} u = \sigma_0\sigma_1 & \quad s_0 \mapsto s_1 \mapsto s_2 \mapsto s_1\uncertain{\text{ ?}},\quad s_3 \mapsto s''_3 \ \text{(à calculer)} \\ v = \sigma_1\sigma_2 & \quad s_0 \mapsto s_0,\quad s_3 \mapsto s_1 \mapsto s_2 \mapsto s'_3 \end{cases} \]
\note{les lettres $u$ et $v$ sont écrites en surcharge d'autres, noircies. Dans la première ligne, le dernier terme du cycle se lit $s_1$ ; on attendrait $s_0$.}

\struck{Notons que $s_0$, $s_1$, $s_2$, $s_3$}

On aura
\begin{align*}
s''_3 = u s_0 + u e_3 &= s_{\uncertain{1}} + \sigma_0(\sigma_1(e_3)) = s_0 + e_1 + \sigma_0\bigl[(1+\alpha)(e_2 - e_1) + e_3\bigr] \\
&= s_0 + e_1 + (1+\alpha)\, e_2 \struck{\ill{}} + e_3 - e_1
\end{align*}
donc
\[ (5) \qquad s''_3 = s_0 + (1+\alpha)\, e_2 \struck{\ill{}} + e_3 = s_3 + (1+\alpha)\, e_2 \]
ou encore
\[ (6) \qquad s''_3 - s_3 = (1+\alpha)(s_2 - s_0) \]

Calculons $u(s'_3)$ utilisant (2)
\[ \struck{u s'_3 = u s_0 + u e_3 = s_1 + \sigma_0\sigma_1\bigl((1+\alpha)(e_2 - e_1) + e_3\bigr)} = s_1 + (1+\alpha)\, u(s_2 - s_1) \]
\note{le $= s_1 + (1+\alpha)u(s_2 - s_1)$ final est lui aussi traversé par le trait.}
\[ u(s'_3) = u(s_3) + (1+\alpha)\underbrace{u(s_2 - s_1)}_{s_0 - s_2 = -e_2} = s''_3 - (1+\alpha)\, e_2 \underset{(5)}{=} s_3 \]

\ill{} \ill{}, on trouve de même par (6)
\[ \struck{(7) \quad u(s_3) \longmapsto s'} \]
\[ \struck{(7)}\ \ u(s''_3) = u(s_3) + (1+\alpha)\underbrace{u(s_2 - s_3)}_{s_0 - s_1 = -e_1} = s''_3 - (1+\alpha)\, e_1 \]
\note{sous la dernière ligne, coupée par le bord de la feuille : « $= s_0 + (1+\alpha)(e_2 - e_1) + e_3$ ».}

\page{85}
donc
\[ (7) \qquad u : s_3 \longmapsto s''_3 \longmapsto s'_3 \longmapsto s_3 \]
Donc on trouve
\[ (8) \qquad u^3 = \mathrm{id} \]
\struck{\ill{}} ($u \neq \mathrm{id}$ bien sûr).

Notons que $s_3$, $s_1$, $s_2$, $s'_3$ sont dans un même plan (celui d'équation $x + y + z = 1$ !) et ce plan est stable par $\sigma_1$ et $\sigma_2$ (donc par $v = \sigma_1\sigma_2$), \struck{ainsi} donc par \add{$v = \sigma_1\sigma_2$}. Un calcul \uncertain{commode} montre que
\[ (9) \qquad v^5 = \mathrm{id} \iff \boxed{\alpha^2 + \alpha - 1 = 0} \]

Donc si $G$ est le groupe
\[ (10) \qquad G = \Bigl\{ \sigma_0, \sigma_1, \sigma_2 \Bigm| \sigma_0^2 = \sigma_1^2 = \sigma_2^2 = 1,\ \sigma_0\sigma_2 = \sigma_2\sigma_{\uncertain{1}},\ (\sigma_0\sigma_1)^3 = 1,\ (\sigma_1\sigma_2)^5 = 1 \Bigr\} \]
\[ = \mathrm{Aut}(\text{icosaèdre comb. épinglé} \ldots) \]
\note{dans la relation de commutation, le dernier indice se lit 1 ; le contexte (page 83) demande $\sigma_0\sigma_2 = \sigma_2\sigma_0$.}
alors sous réserve que $\alpha^2 + \alpha - 1 = 0$, ce groupe \add{\ill{}} $G$ opère \add{\uncertain{effectivement}} sur $\mathbb{E}^3$. Montrons que l'ensemble des transformés de $s_0$, $\mathrm{Dr}(s_0, s_1)$ et $\mathrm{Pl}(s_0, s_1, s_2)$ forme un icosaèdre géométrique ? Cela signifie que sur \uncertain{tout} corps, les seules relations d'incidence entre ces \struck{droites} sommets, droites et plans sont celles

\page{86}
de l'icosaèdre.
\note{le reste de la feuille est blanc ; la question posée à la page 85 n'est pas poursuivie ici.}

\page{87}
\note{la page s'ouvre, comme la page 86, sur la fin de la phrase interrompue au bas de la page 85.}
\struck{de} de l'icosaèdre combinatoire.

Si nous admettons que $G$ est \add{\uncertain{fini}} d'ordre 120, comme le stabilisateur de $s_0$ contient le groupe \add{$Z$} engendré par $v$ (cyclique d'ordre 5) en plus de $\sigma_1$ --- donc le groupe diédral $\{1, \sigma_1\}\cdot Z$ (produit semi-direct), on voit que le nombre de sommets est un diviseur de $120/10 = 12$. Notons que les six points déjà écrits $s_0, s_1, s_2, s_3, s'_3, s''_3$ sont \emph{distincts}, or les coordonnées sont
\[ (11) \qquad \begin{cases}
s_0 = (0, 0, 0) \\
s_1 = (1, 0, 0) \\
s_2 = (0, 1, 0) \\
s_3 = (0, 0, 1) \\
s'_3 = \bigl(-(1+\alpha), 1+\alpha, 1\bigr) = (\alpha', -\alpha', 1) \\
s''_3 = (0, 1+\alpha, 1) = (0, -\alpha', 1)
\end{cases} \]
\note{dans l'expression de droite de $s'_3$, le premier signe est noirci : \uncertain{$\alpha'$}. En marge droite des coordonnées : $\alpha' = -(\alpha+1) = -\alpha^{-1}$ (« l'antiracine » de $\alpha^2 + \alpha - 1 = 0$)  ; le mot entre guillemets est lu avec doute. Des accolades au crayon relient $s_0, s_1$ et $s_2, s_3$.}

\struck{\ill{}} Il est clair que $s_0$ à $s_3$ sont deux à deux distincts, il faut \struck{voir} les comparer à $s'_3$, $s''_3$ et comparer ceux-ci entre eux. Comme la dernière coordonnée de $s'_3$, $s''_3$ \struck{\ill{}} est 1, ils sont $\neq s_0, s_1, s_2$ dont la dernière coordonnée est nulle. Ils \struck{\ill{}} \add{sont} \uncertain{égaux} \add{l'un à l'autre} \uncertain{à} $s_3$ \struck{\ill{}} \add{\uncertain{ssi}} $1 + \alpha = 0$ (i.e.\ $\alpha = -1$), qui est \add{donc} aussi la condition \uncertain{nécessaire} et suffisante pour que $s'_3 = s''_3$. Or \struck{$1+\alpha$} \ill{} si $\alpha = -1$
\note{les deux dernières phrases sont rapides et chargées de surcharges ; leur lecture est incertaine dans le détail. La phrase se poursuit au-delà du bas de la page.}

\page{88}
l'équation $\alpha^2 + \alpha - 1 = 0$ devient $-1 = 0$, qui n'est jamais vérifiée (sur une base non vide !). Cela \struck{signifie} \add{implique} que le nombre de \struck{\ill{}} \add{sommets est} au moins six, donc égal à six ou douze. Si c'était \struck{douze} six, le stabilisateur de \struck{\ill{}} $s_0$ serait un groupe d'ordre 20 contenant le groupe diédral \add{$D_5$} --- or il n'y en a qu'un seul, qui est le produit de $D_5$ par le centre $\mathfrak{Z} = \{1, a\}$ ($a$ = antipodisme). On voit que cela signifierait que $a s_0 = s_0$, ou $a s_3 = s_3$ (cela revient au même par homogénéité). Nous allons donc calculer les antipodiques des 6 points déjà envisagés (qui nous fournissent) --- ce qui nous donnera une liste de 12 pts !
\note{le $\mathfrak{Z}$ est une sorte de 3 bouclé, pour le centre ; « ($a$ = antipodisme) » est écrit sous l'accolade. La parenthèse « (qui nous fournissent) » est lue avec doute.}

On va d'abord déterminer le \emph{centre} de l'icosaèdre, défini (s'il existe) comme l'unique (?) pt fixe sous $G$.

\page{89}
\note{la feuille est vide sauf un calcul isolé dans le quart inférieur droit, relations entre $\alpha$ et son conjugué $\alpha'$.}
\[ \alpha + \alpha' = -1 \qquad \begin{cases} \alpha = -(\alpha' + 1) \\ \alpha' = -(\alpha + 1) \end{cases} \]
\[ \alpha\alpha' = -1 \qquad \begin{cases} \alpha = -1/\alpha' \\ \alpha' = -1/\alpha \end{cases} \]
\begin{gather*}
(1 - \alpha)(1 - \alpha') = 1 \\
(1 + \alpha)\alpha = (1 + \alpha')\alpha' = 1 \qquad (1 + \alpha)(1 + \alpha') = -1 \\
(\alpha + 2)(\alpha' + 2) = 1 \\
(2\alpha - 1)(2\alpha' - 1) = -1
\end{gather*}
\note{$(2\alpha-1)(2\alpha'-1) = 4\alpha\alpha' - 2(\alpha+\alpha') + 1 = -1$ ; la relation est transcrite telle quelle.}

\page{90}
\[ (12) \quad \begin{cases} \sigma_0(x, y, z) = \bigl(1 - (x + y + z), y, z\bigr) \\ \text{pts fixes définis par } 2x + y + z - 1 = 0 \quad \struck{\ill{}} \quad (\text{si } 2 \text{ inv.}) \end{cases} \]
\note{à droite de (12), deux lignes biffées : \struck{\ill{} car. 2, \ill{} l'hyperplan : l'infini\ldots}.}
\[ (13) \quad \begin{cases} \sigma_1(x, y, z) = \bigl(y - (1+\alpha)z,\ x + (1+\alpha)z,\ z\bigr) = (y + \alpha' z,\ x - \alpha' z,\ z) \\ \text{pts fixes définis par } x + (1+\alpha)z = y \quad (\text{vis. en } x, y, z) \\ \phantom{\text{pts fixes définis par }} \struck{x + \alpha' z = y} \quad x - y - \alpha' z = 0 \end{cases} \]
\[ (14) \quad \begin{cases} \sigma_2(x, y, z) = (x, z, y) \\ \text{pts fixes définis par } y = z \end{cases} \]
\note{« vis. » (sans doute « visiblement ») est lu avec doute. Dans (13), les signes devant $\alpha' z$ sont noircis et lus d'après $\alpha' = -(1+\alpha)$.}

\emph{Centre} $o$ de coordonnées
\[ x = \tfrac{1}{2}\alpha', \quad y = \alpha' x = \struck{\ill{}}\ \tfrac{1 - \alpha'}{2}, \quad z = y \qquad \struck{x = \tfrac{\ill{}}{2},\ y = z = -\tfrac{\ill{}}{2\alpha}} \]
\note{la valeur intermédiaire de $y$ est noircie ; « $\frac{1-\alpha'}{2}$ » est écrit au-dessus. La relation « $y = \alpha' x$ » est transcrite telle que lue.}
\[ (15) \qquad \struck{o = \tfrac{1}{2}(\ill{}) = -\tfrac{1}{2}\Bigl(-1, \tfrac{1}{\alpha}, \tfrac{1}{\alpha}\Bigr)} \]
\[ o = \tfrac{1}{2}(\alpha', 1 - \alpha', 1 - \alpha') = \tfrac{1}{2}\bigl(-(1+\alpha), 2 + \alpha, 2 + \alpha\bigr) \]

\emph{NB} L'équation $\alpha^2 + \alpha - 1$ implique $\alpha$ inv.\ et \struck{en fait \ill{}} elle s'écrit
\[ (16) \qquad \boxed{\alpha(\alpha + 1) = 1} \quad \text{ou} \]
\struck{\ill{} $\alpha + 2$ \ill{}} [en fait, $\alpha(\alpha+1)(\alpha-1)(\alpha+2)$ \add{$(2\alpha - 1)$} sont inv.\ \ldots\ \}, et $(\alpha - 2)$ sera inv.\ \uncertain{ssi} 5 est inv.\ \ldots]

Ainsi, en car.\ $\neq 2$ le centre est donné par les coordonnées (15), et en car.\ 2 \add{sans doute} c'est le pt à l'infini dans la direction $\bigl(-1, \frac{1}{\alpha}, \frac{1}{\alpha}\bigr)$ \ldots].
\note{« sera inv.\ ssi 5 est inv. » : $(\alpha-2)(\alpha+3) = \alpha^2+\alpha-6 = -5$ ; lecture probable mais le mot « ssi » est très abrégé.}

\page{91}
Ceci dit, on doit pouvoir \add{sans mal} simplifier ceci.

\emph{Prop.} L'antipodisme \add{$\underline{a}$} est la symétrie par rapport au pt $o$ donné par (15) \add{si 2 inv.}, et dans tous les cas il est donné par la formule
\[ (17) \qquad \underline{a}(x, y, z) = \bigl(\alpha' - x,\ (1 - \alpha') - y,\ (1 - \alpha') - z\bigr) \]
\note{sous (17), une première version biffée et illisible : \struck{\ill{}}. Les premiers termes des trois coordonnées sont écrits en surcharge.}

Si 2 est inv., \add{$\underline{a}$} \struck{\ill{}} a $o$ comme unique pt fixe --- or les $s_i$, \struck{$s_3$} $s'_3$, $s''_3$ sont $\neq o$ (n'étant pas fixes sous $G$). En car.\ \struck{\ill{}} 2, $\underline{a}$ est la translation par le vecteur $(\alpha', \uncertain{1}, \uncertain{1})$, donc n'a pas de pt fixe (à distance finie).
\note{en caractéristique 2, (17) donne la translation par $(\alpha', 1 + \alpha', 1 + \alpha')$ ; les deux dernières composantes sont écrites en surcharge et lues « 1 ».}

\emph{NB} En car.\ 2, \struck{les solutions de} l'équation $\alpha^2 + \alpha - 1 = 0$ sur une ext.\ \add{alg.\ close $K$} de $\mathbb{F}_2$ sont les deux éléments de \struck{\ill{}} $\mathbb{F}_4$ \uncertain{$\smallsetminus \mathbb{F}_2$}. $\widetilde{K} \smallsetminus K$, où $K$ est l'unique ext.\ quadratique de $\mathbb{F}_2$ dans $K$. Si on pose
\[ \mathbb{F}_4 \overset{\text{déf}}{=} \mathbb{F}_2[T]/(T^2 + T + 1) = \text{corps des racines primitives } 3^{\text{ièmes}} \text{ de l'unité sur } \mathbb{F}_2 \]
la donnée d'une $\alpha$ sur $S/\mathbb{F}_2$ \struck{\ill{}} satisfaisant l'équation $\alpha^2 + \alpha - 1 = 0$ équivaut à la donnée d'une $\mathbb{F}_4$-structure sur $S$, i.e.\ d'un morphisme $S \to \operatorname{Spec} \mathbb{F}_4$. Notons que $\frac{1}{\alpha} = \alpha + 1$ est l'\emph{autre} él.\ de $\mathbb{F}_4 \smallsetminus \mathbb{F}_2$.
\note{la phrase sur $\widetilde{K} \smallsetminus K$ est lue telle quelle ; les deux $K$ désignent manifestement la clôture algébrique et l'extension quadratique, que la page note de la même lettre.}

\page{92}
\[ (18) \qquad \begin{cases}
\underline{a}\, s_0 = (\alpha', 1 - \alpha', 1 - \alpha') \\
\underline{a}\, s_1 = (\alpha' - 1, 1 - \alpha', 1 - \alpha') \\
\underline{a}\, s_2 = (\alpha', -\alpha', 1 - \alpha') \\
\underline{a}\, s_3 = (\alpha', 1 - \alpha', -\alpha') \\
\underline{a}\, s'_3 = (0, 1 \ill{}, -\alpha') \\
\underline{a}\, s''_3 = (\alpha', 1 \ill{}, -\alpha')
\end{cases} \]
\note{dans les deux dernières lignes, la deuxième coordonnée est noircie après le « 1 » ; sous $\underline{a}\,s'_3$, une correction également noircie, où l'on devine \uncertain{$2\alpha + 3$}. D'après (11) et (17) on attendrait $\underline{a}\,s'_3 = (0, 1 - 2\alpha', -\alpha')$ et $\underline{a}\,s''_3 = (\alpha', 1, -\alpha')$. Le reste de la feuille (papier à musique) est blanc.}

\page{93}
\note{le haut de la page (titre et tableau (18)) est barré de deux longs traits obliques ; c'est une autre version de (18), avec d'autres valeurs, transcrite ici sous le trait.}
\struck{Écrivons les antipodiques des $s_i$, $s'_3$, $s''_3$}
\[ \struck{(18)} \quad \begin{cases}
\underline{a}\, s_0 = \bigl(1, -\frac{1}{\alpha}, -\frac{1}{\alpha}\bigr) = (1, -1-\alpha, -1-\alpha) = (1, \alpha', \alpha') \\
\underline{a}\, s_1 = \bigl(0, -\frac{1}{\alpha}, -\frac{1}{\alpha}\bigr) = (0, -1-\alpha, -1-\alpha) = (0, \alpha', \alpha') \\
\underline{a}\, s_2 = \bigl(1, -\frac{1}{\alpha} - 1, -\frac{1}{\alpha}\bigr) = (1, -2-\alpha, -1-\alpha) = (1, \alpha' - 1, \alpha') \\
\underline{a}\, s_3 = \bigl(1, -\frac{1}{\alpha}, -\frac{1}{\alpha} - 1\bigr) = (1, \uncertain{-1-\alpha}, -2-\alpha) = (1, \alpha', \alpha' - 1) \\
\underline{a}\, s'_3 = \bigl(2 + \alpha, -2(1+\alpha), -(2+\alpha)\bigr) = (1 - \alpha', -2\alpha', \alpha' - 1) \\
\underline{a}\, s''_3 = \bigl(1, -2(1+\alpha), -(2+\alpha)\bigr) = (1, 2\alpha', \alpha' - 1)
\end{cases} \]

On aimerait prouver que (\add{sur un corps de base quelconque}) \add{trois \uncertain{sommets}} \struck{\ill{}} distincts \add{(ce qui \ill{} \ill{} \ill{} \ill{} \ill{})} ne sont pas alignés, et quatre points distincts ne sont pas dans \struck{des} \add{un} même plan \emph{sauf} s'ils \struck{\ill{}} sont incidents \add{combinatoirement} (à une même face) (ce qui implique que un sommet n'est incident \add{géom.} à une face que s'il \struck{\ill{}} l'est combinatoirement ; et \add{par suite} qu'une arête n'est incidente \add{géom.} à une face que si elle est incidente combinatoirement).
\marginal{« ? » en marge de ce paragraphe, qu'un trait vertical encadre ; plus bas, reliée par une flèche : « c'est en fait \emph{faux} en car.\ 2 ! »}
\note{les ajouts interlinéaires au-dessus des deux premières lignes du paragraphe sont serrés et en partie illisibles.}

Pour le premier point, il y a trois vérifications à faire :

$\bullet$) La droite qui joint $s_0$ à son antipodique $\underline{a}\,s_0$ ne contient aucun autre sommet. \uncertain{On} \uncertain{constate} que les \add{10} vecteurs \struck{\ill{}} \add{$(x, y, z)$} associés aux dix sommets distincts de $s_0$, $\underline{a}\,s_0$, sont \uncertain{non prop.}\ au vecteur \uncertain{de Coxeter} $\underline{a}\,s_0 = (\alpha', 1 - \alpha', 1 - \alpha')$ \struck{$= (\ill{}, -1-\alpha, -1-\alpha)$} ou encore à $\frac{1}{\alpha'}\,\underline{a}\,s_0 = (1, \alpha', \alpha')$
\note{la dernière ligne est serrée contre le bord inférieur ; « de Coxeter » est une lecture très douteuse d'un mot abrégé. La phrase se poursuit en haut de la page 94.}

\page{94}
\emph{Ok} C'est vrai.

$\ast$ b) La droite \emph{joignant} $s_0$ à $s_1$ ne contient pas d'autre sommet, i.e.\ les 10 \struck{vecteurs}-sommets distincts de $s_0$, $s_1$ sont non prop.\ à $s_1 = (1, 0, 0)$. C'est vrai.

c) La droite \emph{joignant} $s_0$ à $s''_3$ ne contient pas d'autre sommet, i.e.\ les 10 vect.-sommets distincts de $s_0$, $s''_3$ sont non prop.\ à $s''_3 = (0, 1 + \alpha, 1) = (0, -\alpha', 1)$. C'est vrai.

Le deuxième pt semble plus fastidieux. Mais si on s'est seulement \uncertain{restreint} au cas de 4 pts \add{distincts} dont trois forment face, on est ramené à ceci :

Aucun des 9 vecteurs-sommets distincts de $s_0, s_1, s_2$ n'est dans le plan qu'ils engendrent, i.e.\ n'a de coordonnée $z$ nulle. Or c'est vrai --- et on est heureux !
\note{le reste de la feuille est blanc.}

\page{95}
Fonctions \struck{\ill{}} quadr.\ \uncertain{invariantes}
\[ f(x, y, z) = ax^2 + by^2 + cz^2 + uyz + vzx + wxy + \lambda x + \mu y + \nu z + d \]
Comme les cstes sont invariantes, OPS $d = 0$. L'invariance par $\sigma_2 : (x, y, z) \mapsto (x, z, y)$ signifie $b = c$, $v = w$, $\mu = \nu$, i.e.
\[ f(x, y, z) = ax^2 + b(y^2 + z^2) + uyz + vx(y + z) + \lambda x + \mu(y + z) \]
Faisant la substitution $\sigma_0 : (x, y, z) \mapsto \bigl(1 - (x + y + z), y, z\bigr)$
\begin{align*}
f \circ \sigma_0(x, y, z) &= a\bigl(1 + x^2 + y^2 + z^2 + 2yz + 2zx + 2xy - 2(x + y + z)\bigr) \\
&\quad + b(y^2 + z^2) + uyz + v\bigl(1 - (x + y + z)\bigr)(y + z) \\
&\quad + \lambda\bigl(1 - (x + y + z)\bigr) + \mu(y + z) \\
&= ax^2 + (b + a - v)(y^2 + z^2) + (u + 2a - 2v)yz \\
&\quad + (2a - v)x(y + z) + (-2a - \lambda)x \\
&\quad + (-2a - \lambda + v + \struck{\mu})(y + z) + (a + \struck{\ill{}} + \lambda)
\end{align*}
\note{les termes sont cochés d'une croix au fur et à mesure du développement ; plusieurs coefficients portent des surcharges (le $b$ de $(b + a - v)$ et le $u$ de $(u + 2a - 2v)$ sont écrits sur d'autres lettres). Le coefficient de $(y+z)$ est lu $-2a - \lambda + v$, le $\mu$ biffé ; le développement complet donnerait $-2a - \lambda + v + \mu$, comme la ligne « $\mu = -2a - \lambda + v + \mu$ » ci-dessous le suppose.}

d'où
\begin{gather*}
\boxed{v = a} \\
2v = 2a \quad \text{ok} \\
2a - v = v \quad \text{ok} \\
-2a - \lambda = \lambda \quad \text{i.e.} \quad \boxed{2(\lambda + a) = 0} \\
\mu = -2a - \lambda + v + \mu \quad \text{i.e.\ (connu)} \quad \boxed{\lambda = -a}
\end{gather*}
\note{la première égalité encadrée a sa lettre de gauche noircie ; d'après les deux lignes suivantes c'est $v = a$. De « $2(\lambda + a) = 0$ » part une flèche vers une expression biffée et encadrée, « \struck{$2\lambda + 2a = 0$} ».}

\struck{$a + v + \lambda = 0 \quad \lambda = 2a \quad \text{ok}$}

\[ f(x, y, z) = ax^2 + b(y^2 + z^2) + uyz + ax(y + z) \add{- ax} \struck{+ \mu(y + z)} \qquad \boxed{\uncertain{2a = 0}} \]
\[ \struck{= ax(x + y + z) + b(y^2 + z^2) + uyz + \mu(y + z)} \]
\[ \struck{= a\bigl[x^2 + x(y + z) - x\bigr]} \]
\[ f(x, y, z) = ax(x + y + z - 1) + b(y^2 + z^2) + uyz + \mu(y + z) \]
\note{l'encadré à droite de l'avant-dernière ligne est lu avec doute ; la dernière ligne garde le terme $\mu(y+z)$ que la précédente biffe.}

\page{96}
\note{toute la page est barrée de deux longs traits obliques.}
Fonction ss-quadratique nulles sur les sommets de l'icosaèdre :
\[ \begin{array}{lcll}
f(s_0) = 0 & \Longrightarrow & d = 0 & \\
f(s_1) = 0 & \Longrightarrow & a + \lambda = 0 & \text{i.e.}\ \lambda = -a \\
f(s_2) = 0 & \Longrightarrow & b + \mu = 0 & \phantom{\text{i.e.}}\ \mu = -b \\
f(s_3) = 0 & \Longrightarrow & c + \nu = 0 & \phantom{\text{i.e.}}\ \nu = -c
\end{array} \]
\[ f(x, y, z) = a(x^2 - x) + b(y^2 - y) + c(z^2 - z) + uyz + vzx + wxy \]
\[ f(s''_3) = b(1+\alpha)\alpha + u(1+\alpha) = 0 \qquad \underline{u = -\alpha b} \]
\begin{align*}
f(\underline{a}\,s_1) &= \uncertain{\alpha}(\alpha+1)(\alpha+2)(b + c) - b(\alpha + 2) = 0 \\
\text{\uncertain{d'où}}\quad & (\alpha+1)(b + c) - b = 0 \quad \text{i.e.} \quad \alpha b + (\alpha+1)c = 0 \\
\text{i.e.}\quad & b = -\Bigl(1 + \frac{1}{\alpha}\Bigr)c = -(2 + \alpha)c
\end{align*}
\note{le mot « ss-quadratique » (sans doute « sous-quadratique » ou « ss.-quadr. ») est lu avec doute. Le premier facteur de $f(\underline{a}\,s_1)$ et le signe devant $b(\alpha+2)$ sont écrits en surcharge.}

\page{97}
\note{la page reprend la forme de $f$ obtenue à la fin de la page 95 et lui applique $\sigma_1$ (formule (13)).}
\begin{align*}
f\sigma_1(x, y, z) &= a(y + \alpha' z)(x + y + z - 1) + b\bigl(x^2 + (\alpha'^2 + 1)z^2 - 2\alpha' xz\bigr) \\
&\quad + u(x - \alpha' z)z + \mu\bigl(x + (1 - \alpha')z\bigr) \\
&= bx^2 + ay^2 + \bigl(a\alpha' + b(1 + \alpha'^2) \add{- u\alpha'}\bigr)z^2 \\
&\quad + (\uncertain{-a\alpha})yz + (a\alpha' - 2b\alpha' + u)zx + axy \\
&\quad + \mu x + (-a)y + \bigl(-a\alpha' + \mu(1 - \alpha')\bigr)z
\end{align*}
\note{les produits développés sont reliés par des accolades et des flèches ; le « $- u\alpha'$ » est inséré au-dessus de la ligne, entouré. Le coefficient de $yz$ est écrit en surcharge.}

d'où $\boxed{b = a}$
\begin{gather*}
a\alpha' + b(1 + \alpha'^2) - u\alpha' = b \\
\text{i.e.} \quad a(\underbrace{\alpha'^2 + \alpha'}_{1} + \struck{\ill{}}\,\uncertain{1}) = u\alpha' \quad \text{i.e.} \quad \boxed{u = -\alpha a} \\
-a\alpha = u \quad \text{ok} \\
a\alpha' - 2b\alpha' + u = a \quad \text{i.e.} \quad a(\uncertain{-}\alpha' - 2\alpha' \struck{\ill{}}\,\underbrace{-\alpha}_{\alpha'} - 1) = 0 \quad \text{ok} \\
a = a \quad \text{ok} \\
\boxed{\mu = -a} \\
-a = \mu \quad \text{ok} \\
-a\alpha' + \mu(\struck{1} - \alpha') = \struck{\mu} \quad \text{i.e.} \quad \struck{a(-\alpha' - a\alpha')}\ -\alpha'(a + \mu) = 0 \quad \text{ok}
\end{gather*}
\note{les signes dans le « i.e. » de la ligne $a\alpha' - 2b\alpha' + u = a$ sont noircis ; le $1$ et le $\mu$ biffés de l'avant-dernière ligne le sont d'un trait fin.}

\[ f(x, y, z) = a\, q(x, y, z) \]
\[ (19) \qquad \boxed{\begin{aligned} q(x, y, z) &= x(x + y + z - 1) + (y^2 + z^2) - \alpha yz - (y + z) \\ &= x^2 + y^2 + z^2 - \alpha yz + zx + xy - (x + y + z) \end{aligned}} \]

Les \struck{fonctions} quadr.\ invariantes sont donc celles de la forme $a\,q(x, y, z) + d$, celles qui sont nulles sur $S$ sont celles de la forme $a\,q(x, y, z)$.

\page{98}
\note{toute la page est barrée d'un long trait oblique. Elle s'ouvre sur « $=$ » : c'est la fin d'un développement de $f\sigma_1$ dont le début n'est pas sur la page 97 telle qu'elle se présente, sans doute une première version écrite avec $\alpha$ au lieu de $\alpha'$.}
\begin{align*}
&= bx^2 + ay^2 + \bigl[a(\alpha + 2) + b(\alpha + 3) + u(1 + \alpha)\bigr]z^2 \\
&\quad + \bigl(-2a(1 + \alpha)\bigr)yz + \bigl[2b(1 + \alpha) + u\bigr]zx \\
&\quad + \mu x + (-a)y + \bigl[a(1 + \alpha) + \mu(2 + \alpha)\bigr]z
\end{align*}
\note{le coefficient de $b$ dans le crochet de $z^2$ est écrit en surcharge : \uncertain{$(\alpha + 3)$} ; celui de $yz$ se lit $-2a(1+\alpha)$ avec un premier chiffre surchargé.}

les conditions deviennent
\[ \begin{cases}
a = b \\
a(2\alpha + 5) + u(1 + \alpha) = \underset{a}{b} \\
-2a(1 + \alpha) = u & \text{déjà vu} \\
2b(1 + \alpha) + u = 0 & \text{déjà vu } (b = a) \\
\mu = -a \\
-a = \mu & \text{déjà vu} \\
a(1 + \alpha) + \mu(2 + \alpha) = \mu \\
\quad \text{i.e.}\ (\mu + a)(1 + \alpha) = 0
\end{cases}
\qquad \text{i.e.} \qquad
\begin{cases}
b = a \\
u = -2a\,\dfrac{\alpha + 2}{1 + \alpha} = -2a\alpha(\alpha + 2) = -2a(\alpha + 1) \\
\mu = -a
\end{cases} \]
\note{dans la fraction, le numérateur « $\alpha + 2$ » est surchargé. Les égalités de la deuxième accolade sont transcrites telles quelles.}

\[ f(x, y, z) = a\, q(x, y, z) \quad \text{où} \]
\[ \boxed{q(x, y, z) = x^2 + y^2 + z^2 \struck{\ill{}} - 2(1 + \alpha)yz - (x + y + z)} \]
\note{la dernière parenthèse commence par une lettre surchargée, lue $x$.}

Ainsi $q$ et le ch.\ $+1$ forment une base des fonctions sous-quadratiques invariantes. Parmi celles-ci, les multiples de $q$ sont celles qui s'annulent sur les sommets de l'icosaèdre.
\note{« ch. » est lu d'après la lettre ; on attend « la cste » (la constante). La lecture est incertaine.}

\page{99}
\note{nouvelle reprise, sur papier à musique, du calcul des quadriques passant par les sommets, sans supposer d'abord l'invariance.}
Quadriques \uncertain{circonscrites}
\[ f(x, y, z) = ax^2 + by^2 + cz^2 + uyz + vzx + wxy + \lambda x + \mu y + \nu z + d \]
\[ \begin{array}{lll|l}
f(s_0) = 0 & \boxed{d = 0} & & f(x, y, z) = \\
f(s_1) = 0 & a + \lambda = 0 & \boxed{\lambda = -a} & a(x^2 - x) + b(y^2 - y) + c(z^2 - z) \\
f(s_2) = 0 & b + \mu = 0 & \boxed{\mu = -b} & \quad + uyz + vzx + wxy \\
f(s_3) = 0 & c + \nu = 0 & \boxed{\nu = -c} &
\end{array} \]
\[ f(s''_3) = 0 \qquad b(\alpha'^2 + \alpha') \struck{\ill{}} + u\alpha' = 0 \qquad \boxed{u = \struck{\ill{}} -\alpha b} \]
\[ u = -(\alpha' + 1)b = -\alpha b \]
\note{« circonscrites » est lu avec doute. Le signe entre $b(\alpha'^2+\alpha')$ et $u\alpha'$ est une tache.}

\[ \struck{f = a(x^2 - x) + b\bigl(y^2 - \ill{}\bigr) + c(z^2 - z)} \]
\[ f(x, y, z) = a(x^2 - x) + by(y - \alpha z - 1) + c(z^2 - z) + \struck{\ill{}} \add{x(vz + wy)} \]
\[ f(\struck{\ill{}}\,s'_3) = 0 \qquad b\bigl(\underbrace{1 - \underbrace{\alpha(-\alpha')}_{1} - 1}_{-1}\bigr) + c(\underbrace{\alpha'^2 + \alpha'}_{1}) = 0 \qquad \boxed{b = c} \]
\[ f(x, y, z) = a(x^2 - x) + b\bigl(\struck{y(y - \alpha z - 1) + z(z - 1)}\ \add{y^2 + z^2 - \alpha yz - (y + z)}\bigr) + x(vz + wy) \]

\begin{align*}
f(s'_3) = 0 \qquad & a(\alpha'^2 - \alpha') + b\bigl(-\alpha'(\underbrace{-\alpha' - \alpha - 1}_{0}) + 0\bigr) \\
& \struck{+ \alpha'(}\ v\alpha' - \alpha'^2 w = 0 \\
& a(\alpha' - 1) + v - \alpha' w = 0 \qquad \boxed{v = \alpha' w - (\alpha' - 1)a}
\end{align*}

\begin{align*}
f(\underline{a}\,s''_3) = 0 \qquad & a(\alpha'^2 - \alpha') + b\Bigl(\underbrace{\struck{1} + \underbrace{\alpha\alpha'}_{-1} \struck{- 1} + \underbrace{(-\alpha')(-\alpha' - 1)}_{\alpha'^2 + \alpha' = 1}}_{0}\Bigr) \\
& + (-\alpha'^2)v + \alpha' w = 0 \\
& (\alpha' - 1)a \struck{\ill{}} - \alpha' v + w = 0 \qquad \text{ou encore, \uncertain{tout} \uncertain{court}} \quad \boxed{w = \alpha' v - (\alpha' - 1)a}
\end{align*}
\note{sous l'accolade du « $0$ », une ligne biffée illisible. Le « $-\alpha' v + w$ » est écrit en surcharge d'une première version biffée, où l'on devine \uncertain{$\alpha(\alpha' - 1)$}.}

\struck{Compte tenu des ci-dessus \ldots}
\[ \struck{(\alpha' - 1)a - 2\alpha b + (1 - \alpha'^2)w + (\alpha'^2 - \alpha')a = 0} \]
\[ \struck{(\alpha'^2 - 1)a - 2\alpha b + \underbrace{(1 - \alpha'^2)}_{\alpha'}w = 0} \]
\note{les trois dernières lignes sont barrées de traits obliques ; sous $(\alpha'^2 - 1)$ il note « $-\alpha'$ ». Le calcul s'interrompt ici ; la page 100 traite d'autre chose, et la suite éventuelle de ce calcul n'est pas dans ce lot.}

\page{100}
\note{feuille d'un tout autre sujet (topologie générale : prolongement d'homotopies), au crayon puis à l'encre bleue, barrée de deux longs traits obliques. Sans lien visible avec l'icosaèdre ; elle est transcrite parce qu'elle est de sa main et mathématique.}
$f$ constante dans \uncertain{un} voisinage \add{fermé} \struck{\ill{}} de $x$

$\alpha$ fonction \add{$X \to I$} égale à $1$ en $x$, à $0$ dans $\complement \overset{\circ}{F}$ \quad $\dfrac{\complement\overset{\circ}{F}}{F'}$

\struck{$c(t$}
\[ h : X \times I \longrightarrow Y \]
\[ h \mid F' \times I = (f \mid F') \circ \mathrm{pr}_1 \]
\[ h \mid F \times I = \struck{\ill{}}\bigl[(b, t) \longmapsto c\bigl(t\varphi(b)\bigr)\bigr] \]
\[ f(x) = c(0) \]
\note{les lignes qui précèdent sont au crayon ; la suite est à l'encre. Le « $c$ » de $c(t\varphi(b))$ porte un point au-dessus, lu comme une lettre simple.}

$C \supset C'$ \qquad $f_t$

$X \supset F \supset \overset{\circ}{F}$
\[ f_t(x) = x \quad \text{si}\ x \in \complement\overset{\circ}{F} \]
\[ \struck{g_t(x) = x} \]
\[ F \times I \longrightarrow F \]
\[ \begin{cases} g_t(x) = x & x \in B = F - \overset{\circ}{F} \\ g_{\struck{t}}(s) = s \\ g_1(s) \end{cases} \]

$X = F' \cup F$, \quad $F \cap F' = B$, \quad $x \in F - B$

a) $\exists$ fonction \add{\uncertain{numérique}} continue \uncertain{sur} $F$ égale à $1$ sur $x$, à $0$ sur $B$ (\uncertain{p.\ ex.}\ si $F$ \uncertain{compact} \uncertain{régulier}).
\note{la dernière parenthèse est très abrégée ; « compact régulier » est une lecture douteuse. Le mot inséré au-dessus de « continue » est lu « numérique » avec doute.}

\end{document}
